%=============================================================================
% Section 12: Conclusion
%=============================================================================
\section{Conclusion}

Intelligent behavior alone does not enable AI systems to influence the world. A structured interface between intelligence and action is required. This paper introduced the concept of the AI Contact Layer, a conceptual framework that describes how AI systems interact with external capabilities and environments.

The central claim of this paper is that the Intelligence-Action separation is not merely an architectural convenience, but a fundamental design principle for AI-native systems.

The contact layer provides:
\begin{enumerate}[label=\arabic*.]
\item \textbf{Conceptual Clarity:} A clear distinction between cognition and action
\item \textbf{Architectural Foundation:} A structured approach to implementing action interfaces
\item \textbf{Unified Perspective:} Integration of execution infrastructure with broader system design
\item \textbf{Practical Guidance:} Principles for designing capability-mediated AI systems
\end{enumerate}

The contact layer serves as the bridge between internal reasoning and real-world action, enabling the emergence of AI-native systems capable of meaningful interaction with both digital and physical environments.

As AI systems evolve from information processing to environment interaction, the contact layer will play an increasingly central role in AI architecture. By defining this layer clearly, we provide the conceptual foundation for the next generation of intelligent software infrastructure.
